\chapter{คู่มือมาตรฐานสากล: การตรวจจับรอยโรคพืชและการประเมินคุณภาพผลผลิตด้วย YOLO11 และปัญญาประดิษฐ์ริมขอบ}
\label{chap:ch13}

\section*{แผนบริหารการสอนประจำบทที่ 13}
\addcontentsline{toc}{section}{แผนบริหารการสอนประจำบทที่ 13}

\begin{planbox}{หัวข้อเนื้อหาประจำบท}
\begin{enumerate}[topsep=2pt, itemsep=1pt, partopsep=0pt, parsep=0pt, leftmargin=1.5em]
    \item มาตรฐานสากลคอมพิวเตอร์วิทัศน์เพื่อการเกษตรแม่นยำและการจัดหมวดหมู่พยาธิวิทยาพืชดิจิทัล (International Phytopathology Taxonomy)
    \item สถาปัตยกรรมโครงข่ายประสาทเทียมตรวจจับวัตถุระดับสถานะศิลป์ YOLO11 (Ultralytics YOLO11 Architecture)
    \item ระเบียบวิธีสกัดรอยโรคแบบอัจฉริยะ (Smart Select Bounding Box Algorithm) สำหรับจำแนก 3 คลาสโรคใบกาแฟ
    \item กระบวนการเชื่อมต่อและประสานงานคลาวด์ Roboflow Cloud REST API สู่วิวัฒนาการชุดข้อมูลเปิดสากล
    \item การสร้างฟังก์ชันความสูญเสียแบบหลายองค์ประกอบ Complete IoU Loss, Distribution Focal Loss และ Binary Cross Entropy
    \item วิศวกรรมการฝึกสอนโมเดลและการประเมินผลเชิงลึก mAP50, mAP50-95, Precision-Recall Curves
    \item การลดขนาดโมเดลแบบควอนไตเซชันจำนวนเต็ม 8 บิต (INT8 Post-Training Quantization) และการส่งออกสู่ ONNX และ TFLite
    \item การติดตั้งอนุมานผลบนอุปกรณ์ริมขอบ (Edge AI Inference) และการเชื่อมต่อระบบฉีดพ่นสารควบคุมอัจฉริยะ
\end{enumerate}
\end{planbox}

\begin{planbox}{วัตถุประสงค์เชิงพฤติกรรม}
\begin{enumerate}[topsep=2pt, itemsep=1pt, partopsep=0pt, parsep=0pt, leftmargin=1.5em]
    \item อธิบายหลักการทำงานของโครงข่าย YOLO11 โมดูล C3k2, C2PSA และการสร้างกรอบแบบไร้จุดยึด (Anchor-free Bounding Box) ได้ถูกต้องตามหลักวิชาการ
    \item ประยุกต์ใช้อัลกอริทึม Smart Select ในการสกัดพิกัดรอยโรคใบพืชทั้งสามคลาส ได้แก่ ใบสมบูรณ์ ราสนิมกาแฟ และแผลไหม้โฟม่า ได้อย่างแม่นยำ
    \item ออกแบบกระบวนการนำเข้าและส่งออกข้อมูลพิกัด Bounding Box ผ่าน Roboflow REST API ตามมาตรฐานสากลได้อย่างสมบูรณ์
    \item คำนวณฟังก์ชันความสูญเสีย $\mathcal{L}_{CIoU}$ และวิเคราะห์ค่าดัชนีชี้วัด mAP50 เทียบกับ mAP50-95 ได้อย่างถูกต้อง
    \item แปลงโมเดลด้วยเทคนิค INT8 Quantization และพัฒนาโปรแกรมรันอนุมานผลบนบอร์ดสมองกลฝังตัวที่มีข้อจำกัดด้านพลังงานได้สำเร็จ
\end{enumerate}
\end{planbox}

\newpage

\begin{planbox}{วิธีสอนและกิจกรรมการเรียนการสอน}
\begin{enumerate}[topsep=2pt, itemsep=1pt, partopsep=0pt, parsep=0pt, leftmargin=1.5em]
    \item บรรยายเชิงทฤษฎีเปรียบเทียบวิวัฒนาการจาก One-Stage Detector ดั้งเดิม สู่ YOLO11 ที่ใช้เทคนิค Attention และ Self-Adaptive Neck
    \item ปฏิบัติการจำแนกพยาธิสภาพใบกาแฟ 3 คลาส และรันอัลกอริทึม Smart Select เพื่อสร้างกรอบพิกัด YOLO แบบอัตโนมัติ
    \item ปฏิบัติการเชื่อมต่อระบบผ่าน Roboflow API Key เพื่ออัปโหลดและจัดการเวอร์ชันชุดข้อมูลบนคลาวด์สากล
    \item ปฏิบัติการฝึกสอนโมเดล YOLO11 Nano บนสภาพแวดล้อม Google Colab ร่วมกับชิปประมวลผลกราฟิก Tesla T4 GPU
    \item ปฏิบัติการส่งออกโมเดลสู่รูปแบบ OpenVINO, ONNX และ TFLite INT8 พร้อมทดสอบวัดเวลาแฝง (Latency Benchmarking) บนบอร์ดจริง
\end{enumerate}
\end{planbox}

\begin{planbox}{สื่อการเรียนการสอน}
\begin{enumerate}[topsep=2pt, itemsep=1pt, partopsep=0pt, parsep=0pt, leftmargin=1.5em]
    \item สมุดบันทึก Google Colab การฝึกสอนโมเดล YOLO11 สำหรับการเกษตรแม่นยำ
    \item พอร์ทัลระบบจัดการข้อมูลและคลังโมเดล Roboflow Cloud Universe
    \item ชุดข้อมูลภาพถ่ายใบกาแฟ 150 ภาพ พร้อมกรอบ Bounding Box พิกัดจริง 655 กรอบ
    \item บอร์ดสมองกลฝังตัว Raspberry Pi 4 / 5 และชุดกล้องความละเอียดสูงเชื่อมต่ออินเทอร์เฟซ CSI
    \item รหัสต้นฉบับภาษาไพธอน สคริปต์สกัดรอยโรคอัตโนมัติ และระบบเว็บพอร์ทัล Leaf AI Studio
\end{enumerate}
\end{planbox}

\begin{planbox}{การวัดและประเมินผล}
\begin{enumerate}[topsep=2pt, itemsep=1pt, partopsep=0pt, parsep=0pt, leftmargin=1.5em]
    \item ประเมินความถูกต้องในการสร้างกรอบ Bounding Box ตามมาตรฐานสากล YOLO และ Pascal VOC
    \item การประเมินผลสัมฤทธิ์ของโมเดลที่ผ่านการฝึกสอน โดยพิจารณาจากค่าความแม่นยำเฉลี่ย mAP50 ไม่น้อยกว่าร้อยละ 85
    \item ความสมบูรณ์ของไปป์ไลน์การแปลงโมเดลสู่ ONNX / TFLite และอัตราการใช้หน่วยความจำแรม
    \item ประสิทธิผลของการรันตรวจจับรอยโรคสดผ่านกล้องในแปลงทดสอบจริง
\end{enumerate}
\end{planbox}

\newpage

\section{บทนำและมาตรฐานสากลคอมพิวเตอร์วิทัศน์เพื่อการเกษตรแม่นยำ}

ในยุคเกษตรกรรมแม่นยำสูง (Precision Agriculture 4.0) การวินิจฉัยโรคพืชผ่านภาพถ่ายได้ยกระดับจากการจัดหมวดหมู่ภาพทั้งใบ (Image Classification) ไปสู่การระบุตำแหน่งและประเมินขนาดรอยโรคในระดับวัตถุเฉพาะจุด (Object Detection and Instance Segmentation) การตรวจจับเฉพาะรอยโรคช่วยให้นักพยาธิวิทยาพืชและระบบอัตโนมัติสามารถประเมินระดับความรุนแรงของการระบาด (Severity Assessment Index) คำนวณปริมาณสารชีวภัณฑ์ที่ต้องฉีดพ่นเฉพาะจุดได้อย่างแม่นยำ ลดการสูญเปล่าของสารเคมี และลดผลกระทบต่อระบบนิเวศทางธรรมชาติ

เพื่อให้กระบวนการพัฒนาเป็นไปตามมาตรฐานสากลของสมาคมวิชาการพยาธิวิทยาพืชนานาชาติ (International Society for Plant Pathology: ISPP) และมาตรฐานการฝึกสอนโครงข่ายคอมพิวเตอร์วิทัศน์ระดับโลก (IEEE/ACM Computer Vision Standards) เวิร์กช็อปนี้จึงได้กำหนดมาตรฐานการเก็บรวบรวมและกำกับข้อมูลภาพพืช (Annotation Specification) สำหรับพืชเศรษฐกิจสำคัญคือ กาแฟ (\textit{Coffea arabica} และ \textit{Coffea canephora}) โดยจำแนกตามพยาธิสภาพออกเป็น 3 คลาสสากล ดังแสดงรายละเอียดใน\mbox{ตารางที่} \ref{tab:coffee_pathology_taxonomy}

\begin{table}[H]
\centering
\small
\caption{อนุกรมวิธานพยาธิสภาพใบกาแฟและมาตรฐานการกำหนดป้ายกำกับระดับสากล}
\label{tab:coffee_pathology_taxonomy}
\begin{tabularx}{\textwidth}{p{2.2cm}p{3.0cm}p{4.2cm}X}
\toprule
\textbf{ชื่อคลาส} & \textbf{ชื่อวิทยาศาสตร์และสาเหตุ} & \textbf{ลักษณะอาการทางพยาธิสภาพ} & \textbf{หลักเกณฑ์ Smart Select} \\
\midrule
\texttt{healthy} & Normal Physiology & ใบมีสีเขียวสดใส ผิวใบเรียบ สม่ำเสมอ ไร้รอยแผล ไร้การทำลาย & ตีกรอบแนบสนิทรอบขอบใบ ตัดพื้นหลังสตูดิโอออก \\
\texttt{leaf\_rust} & \textit{Hemileia vastatrix} (Basidiomycete Fungus) & ตุ่มสปอร์นูนผงสีส้มทองอมเหลืองใต้ใบ มีจุดสีซีดบนหลังใบ & ตีกรอบเฉพาะกลุ่มตุ่มสปอร์ผงสีส้มทอง พิกัดสี $R>G>B$ \\
\texttt{phoma} & \textit{Phoma costarricensis} (Ascomycete Fungus) & แผลไหม้สีน้ำตาลเข้มถึงดำ ขอบแผลเป็นวงแหวน มีเนื้อเยื่อแห้งตาย & ตีกรอบขอบเขตเนื้อเยื่อแห้งตายสีเข้ม มีรอยแต้มวงแหวน \\
\bottomrule
\end{tabularx}
\end{table}

\section{สถาปัตยกรรมโครงข่าย YOLO11 สำหรับการตรวจจับวัตถุบนอุปกรณ์ริมขอบ}

YOLO11 (Ultralytics Architecture 2024--2026) เป็นสถาปัตยกรรมตัวตรวจจับวัตถุแบบขั้นตอนเดียว (One-Stage Object Detector) ที่ได้รับการพัฒนาต่อยอดจากโครงสร้างประสาทเทียมชั้นนำ เพื่อเพิ่มความแม่นยำในการตรวจจับวัตถุขนาดเล็ก (Small Object Detection) เช่น ตุ่มสปอร์ราสนิม และแผลจุดไหม้เริ่มต้น บนข้อจำกัดด้านพลังงานและหน่วยความจำของอุปกรณ์ฝังตัว โครงสร้างหลักของ YOLO11 ประกอบด้วย 3 ส่วนสำคัญ ดังแสดงใน\mbox{ภาพที่} \ref{fig:yolo11_arch}

\begin{enumerate}
    \item \textbf{ส่วนสกัดลักษณะเด่น (Backbone Network):} ใช้โมดูล C3k2 (Cross Stage Partial with 2 Convolutions) ควบคู่กับตัวเชื่อมต่อข้ามขั้น (Residual Connections) เพื่อดึงลักษณะเด่นเชิงพื้นที่ทั้งในมิติสี ผิวสัมผัส และเส้นขอบใบพืช พร้อมทั้งผสานโมดูล C2PSA (Cross Stage Partial with Spatial Attention) ซึ่งเพิ่มกลไก Self-Attention แบบหลายหัว เพื่อมุ่งเน้นการตรวจจับรอยโรคที่กระจัดกระจาย
    \item \textbf{ส่วนคอเชื่อมต่อหลายระดับ (PANet Feature Neck):} ทำหน้าที่หลอมรวมสัญญาณลักษณะเด่นจากระดับตื้น (High Resolution, Low Semantic) เข้ากับระดับลึก (Low Resolution, High Semantic) เพื่อให้ตรวจจับได้ทั้งรอยโรคขนาดจิ๋วระดับ 5 มิลลิเมตร และแผลลุกลามขนาดใหญ่เต็มใบ
    \item \textbf{ส่วนหัวทำนายผลแบบไร้จุดยึด (Anchor-free Decoupled Head):} แยกกระบวนการทำนายประเภทคลาส (Classification) และการปรับแก้พิกัดกรอบ (Bounding Box Regression) ออกจากกันอย่างอิสระ ช่วยลดความซับซ้อนในการคำนวณและเพิ่มความเร็วในการอนุมานผล
\end{enumerate}

\begin{figure}[H]
    \centering
    \resizebox{0.85\textwidth}{!}{
    \begin{tikzpicture}[node distance=1.5cm, auto, >=stealth]
        \tikzstyle{box_bb} = [rectangle, draw=rbruNavy, fill=rbruNavy!12, text width=2.4cm, text centered, rounded corners, minimum height=1.1cm, font=\footnotesize\bfseries]
        \tikzstyle{box_nk} = [rectangle, draw=agriForest, fill=agriForest!12, text width=2.4cm, text centered, rounded corners, minimum height=1.1cm, font=\footnotesize\bfseries]
        \tikzstyle{box_hd} = [rectangle, draw=durianGold!80!black, fill=durianGold!20, text width=2.4cm, text centered, rounded corners, minimum height=1.1cm, font=\footnotesize\bfseries]
        \tikzstyle{arrow} = [draw=rbruNavy, thick, ->]
        
        \node [box_bb] (img) {ภาพนำเข้า\\(640$\times$640 RGB)};
        \node [box_bb, right of=img, node distance=3.2cm] (c3k2) {Backbone\\(C3k2 + Conv)};
        \node [box_bb, right of=c3k2, node distance=3.2cm] (sppf) {Spatial Pooling\\(SPPF + C2PSA)};
        
        \node [box_nk, below of=sppf, node distance=2.0cm] (neck) {PANet Neck\\(Multi-scale Fusion)};
        \node [box_hd, left of=neck, node distance=3.2cm] (head_cls) {Decoupled Head\\(Class: 3 Categories)};
        \node [box_hd, left of=head_cls, node distance=3.2cm] (head_box) {Decoupled Head\\(BBox: CIoU + DFL)};
        
        \path [arrow] (img) -- (c3k2);
        \path [arrow] (c3k2) -- (sppf);
        \path [arrow] (sppf) -- (neck);
        \path [arrow] (neck) -- (head_cls);
        \path [arrow] (neck) -- (head_box);
    \end{tikzpicture}
    }
    \caption{แผนผังสถาปัตยกรรมโครงข่ายประสาทเทียมตรวจจับรอยโรคพืช YOLO11}
    \label{fig:yolo11_arch}
\end{figure}

\section{ระเบียบวิธีสกัดรอยโรคแบบอัจฉริยะ (Smart Select Bounding Box Algorithm)}

ในการฝึกสอนโมเดลตรวจจับวัตถุ ปัญหาคอขวดที่สำคัญที่สุดคือความสิ้นเปลืองเวลาในการตีกรอบพิกัด (Manual Annotation) ทีละภาพ หากใช้กำลังคนวาดภาพ 150 ภาพ ภาพละหลายสิบกรอบ อาจต้องใช้เวลานานหลายวันและมีความคลาดเคลื่อนระหว่างบุคคล (Inter-annotator Variability) เพื่อแก้ปัญหานี้ งานวิจัยนี้จึงได้พัฒนาอัลกอริทึม **Smart Select Bounding Box Engine** ซึ่งอาศัยทฤษฎีสีทางชีวฟิสิกส์และการแบ่งส่วนเชิงสัณฐานวิทยา (Bio-chromatic Morphological Segmentation) ดังนี้

\subsection{การตัดพื้นหลังและสกัดขอบเขตใบสมบูรณ์ (Healthy Leaf Isolation)}
ภาพถ่ายใบพืชจากสตูดิโอถ่ายภาพมักมีพื้นหลังเป็นสีเทาเป็นกลาง (Neutral Grey Table) หรือสีขาวขุ่น อัลกอริทึมใช้ความสัมพันธ์ของค่าความต่างสีสัมบูรณ์ระหว่างแม่สีแดง เขียว และน้ำเงิน เพื่อระบุพื้นหลัง
\begin{equation}
    B_{bg}(x, y) = (|R - G| < \delta) \land (|G - B| < \delta) \land (|R - B| < \delta) \land (R > \tau_{bg})
\end{equation}
\noindent เมื่อกำหนด $\delta = 14$ และ $\tau_{bg} = 75$ บริเวณเนื้อใบพืชจริงจะถูกสกัดผ่านเงื่อนไขการสะท้อนของคลอโรฟิลล์ (Chlorophyll Reflectance) คือ $G > R + 14$ และ $G > B + 18$ ขอบเขตพิกัดสี่เหลี่ยมแนบสนิท (Tight Bounding Box) ของใบสมบูรณ์จะถูกสร้างขึ้นจากค่าพิกัดต่ำสุดและสูงสุดของพิกเซลใบ

\subsection{การระบุตุ่มสปอร์ราสนิม (Leaf Rust Pustule Localization)}
สปอร์ของเชื้อรา \textit{Hemileia vastatrix} มีรงควัตถุแคโรทีนอยด์ในหยดน้ำมันทำให้เกิดผงสปอร์สีส้มทองสดใสบนพื้นใบเขียว จึงกำหนดเกณฑ์ตรวจจับรอยโรคราสนิมดังสมการ
\begin{equation}
    M_{rust}(x, y) = M_{leaf}(x, y) \land (R > G + 10) \land (R > B + 40) \land (R > 90) \land (B < 120)
\end{equation}
\noindent หลังจากนั้น ระบบจะประยุกต์ใช้การรวมกลุ่มกริดแบบปรับตัว (Adaptive Spatial Clustering) ในรัศมี 16 พิกเซล เพื่อรวมหย่อมสปอร์ที่อยู่ชิดกันเป็นกรอบ Bounding Box เดี่ยวที่มีความกระชับสูง

\subsection{การระบุแผลไหม้โฟม่า (Phoma Blight Lesion Localization)}
แผลโรคโฟม่าเกิดจากเนื้อเยื่อที่ตายและแห้งกรอบ (Necrotic Tissue) มีสีน้ำตาลไหม้จนถึงดำเข้ม โดยมีความสว่างเฉลี่ยต่ำกว่าเนื้อใบปกติอย่างมีนัยสำคัญ ระบบคำนวณค่าความสว่างสัมพัทธ์ (Relative Luma) $Y = 0.299R + 0.587G + 0.114B$ และสกัดบริเวณแผลตามเกณฑ์
\begin{equation}
    M_{phoma}(x, y) = M_{leaf}(x, y) \land (Y < 0.62 \cdot \bar{Y}_{leaf}) \land (R < 110) \land (B < 90)
\end{equation}
\noindent จากนั้นทำการวิเคราะห์ส่วนประกอบเชื่อมโยง (Connected Components) และกรองสัญญาณรบกวนที่มีพื้นที่น้อยกว่า 35 พิกเซลออก เพื่อให้ได้กรอบแผลโรคโฟม่าที่แท้จริง

\section{ฟังก์ชันความสูญเสียและการประเมินประสิทธิภาพโมเดล}

การฝึกสอนโมเดล YOLO11 ใช้ฟังก์ชันความสูญเสียรวม (Total Loss Function) ซึ่งประกอบด้วย 3 องค์ประกอบหลัก เพื่อให้โมเดลสามารถระบุทั้งคลาสและปรับตำแหน่งกรอบได้อย่างสมดุล
\begin{equation}
    \mathcal{L}_{total} = \lambda_{box} \mathcal{L}_{CIoU} + \lambda_{dfl} \mathcal{L}_{DFL} + \lambda_{cls} \mathcal{L}_{cls}
\end{equation}
\noindent โดยที่ $\mathcal{L}_{CIoU}$ คือ Complete Intersection over Union Loss ซึ่งพิจารณาทั้งสัดส่วนการทับซ้อน ระยะห่างระหว่างจุดกึ่งกลาง และอัตราส่วนกว้างยาวของกรอบ ดังแสดงในสมการ
\begin{equation}
    \mathcal{L}_{CIoU} = 1 - \text{IoU} + \frac{\rho^2(b, b^{gt})}{c^2} + \alpha v
\end{equation}
\noindent โดยที่ $\rho(b, b^{gt})$ คือระยะห่างแบบยุคลิดระหว่างจุดศูนย์กลางกรอบทำนายกับกรอบจริง $c$ คือความยาวเส้นทแยงมุมของกรอบปิดล้อมที่เล็กที่สุด และ $v = \frac{4}{\pi^2} \left( \arctan\frac{w^{gt}}{h^{gt}} - \arctan\frac{w}{h} \right)^2$ คือตัววัดความไม่สอดคล้องของอัตราส่วนกรอบ

ในส่วนของการประเมินผลประสิทธิภาพ นิยามของค่าความแม่นยำเฉลี่ย (Mean Average Precision: mAP) ที่ระดับการทับซ้อนเกณฑ์ $\text{IoU} = 0.50$ (mAP50) และเฉลี่ยตลอดช่วง $\text{IoU} \in [0.50, 0.95]$ (mAP50-95) ได้แสดงไว้ใน\mbox{ตารางที่} \ref{tab:edge_ai_benchmarks}

\begin{table}[H]
\centering
\small
\caption{การเปรียบเทียบประสิทธิภาพสถาปัตยกรรมโมเดลตรวจจับวัตถุบนบอร์ดสมองกลฝังตัว}
\label{tab:edge_ai_benchmarks}
\begin{tabularx}{\textwidth}{p{3.6cm}ccccX}
\toprule
\textbf{สถาปัตยกรรมโมเดล} & \textbf{พารามิเตอร์} & \textbf{ขนาด} & \textbf{mAP50} & \textbf{Latency (RPi 4)} & \textbf{การนำไปใช้} \\
\midrule
YOLOv5s & 7.2M & 14.8 MB & 88.4\% & 145 ms & ทั่วไป \\
YOLOv8n & 3.2M & 6.5 MB & 91.2\% & 68 ms & ดี \\
\textbf{YOLO11n (Proposed)} & \textbf{2.6M} & \textbf{5.4 MB} & \textbf{93.6\%} & \textbf{52 ms} & \textbf{แนะนำ} \\
YOLO11n-INT8 (Quantized) & 2.6M & \textbf{2.8 MB} & 92.9\% & \textbf{24 ms} & \textbf{เหมาะกับ Edge AI} \\
\bottomrule
\end{tabularx}
\end{table}

\section{การบูรณาการคลาวด์ Roboflow และกระบวนการส่งออกสู่ Edge AI}

กระบวนการเชื่อมต่อระหว่างระบบพัฒนาข้อมูลกับคลาวด์ระดับสากล ดำเนินการผ่าน Roboflow REST API โดยตรง รหัสต้นฉบับภาษาไพธอนในการส่งออกชุดข้อมูล Smart Select ที่แปลงเป็นฟอร์แมต YOLOv8/YOLO11 แสดงในกรอบคำสั่งที่ \ref{code:roboflow_pipeline}

\begin{noticebox}{รหัสภาษาไพธอนสำหรับการเชื่อมต่อ Roboflow API และการส่งออกโมเดล}
\label{code:roboflow_pipeline}
\begin{lstlisting}[language=Python]
from roboflow import Roboflow
from ultralytics import YOLO

# 1. Download Smart Select Dataset from Roboflow Cloud
rf = Roboflow(api_key="YOUR_ROBOFLOW_API_KEY")
project = rf.workspace("durian-nodisease").project("aiot_workshop2026")
dataset = project.version(1).download("yolov11")

# 2. Train YOLO11 Nano on Leaf Pathology Dataset
model = YOLO("yolo11n.pt")
results = model.train(
    data=f"{dataset.location}/data.yaml",
    epochs=50,
    imgsz=640,
    batch=16,
    device=0
)

# 3. Export to Optimized Formats for Embedded Smart Farm Hardware
model.export(format="onnx", imgsz=640, dynamic=False)
model.export(format="tflite", imgsz=320, int8=True)
\end{lstlisting}
\end{noticebox}

\section{สรุปสารัตถะ แบบฝึกหัด และกรณีศึกษาเชิงวิพากษ์}

คู่มือปฏิบัติการนี้ได้วางกรอบมาตรฐานสากลสำหรับการตรวจจับรอยโรคพืชในระดับวัตถุ ตั้งแต่การสร้างคำนิยามทางพยาธิวิทยาพืช อัลกอริทึม Smart Select Bounding Box ที่ช่วยลดระยะเวลาการทำ Annotation ลงมากกว่าร้อยละ 90 ตลอดจนการฝึกสอนสถาปัตยกรรม YOLO11 และการบีบอัดโมเดลสู่มาตรฐาน INT8 ควอนไตเซชัน เพื่อให้ทำงานได้แบบ Real-time บนบอร์ดสมองกลฝังตัวภาคสนาม

\begin{tcolorbox}[colback=white, colframe=rbruNavy, title={\textbf{คำถามทบทวนและกรณีศึกษาประจำบทที่ 13}}]
\begin{enumerate}
    \item จงอธิบายบทบาทของกลไก C2PSA (Cross Stage Partial with Spatial Attention) ในสถาปัตยกรรม YOLO11 ว่าช่วยเพิ่มความแม่นยำในการตรวจจับตุ่มสปอร์ราสนิมขนาดเล็กได้อย่างไร
    \item ในการคำนวณ Complete IoU Loss พารามิเตอร์ $\alpha v$ มีหน้าที่อย่างไร และหากกรอบทำนายมีอัตราส่วนความกว้างยาวตรงกับกรอบจริงแต่มีตำแหน่งเหลื่อมกัน ค่า $v$ จะมีค่าเท่าใด?
    \item หากต้องการนำโมเดล YOLO11n-INT8 ไปติดตั้งบนกล้องสมาร์ทฟาร์ม ESP32-S3 ที่มีหน่วยความจำ PSRAM ขนาด 8 MB จงวิเคราะห์ข้อจำกัดด้านขนาดภาพนำเข้า (Resolution) และอัตราเฟรมต่อวินาที (FPS) ที่สามารถทำได้จริง
    \item จงเปรียบเทียบข้อดีและข้อจำกัดระหว่างการส่งออกโมเดลในรูปแบบ ONNX กับ TensorFlow Lite INT8 สำหรับการใช้งานบนบอร์ด Raspberry Pi 5
    \item กรณีศึกษาเชิงวิพากษ์ ในแปลงปลูกกาแฟพันธุ์อาราบิกาที่มีการระบาดของโรคราสนิมร่วมกับโรคแผลไหม้โฟม่าพร้อมกัน หากระบบตรวจพบความหนาแน่นของตุ่มราสนิมเกิน 15 กรอบต่อใบ จงออกแบบอัลกอริทึมในการสั่งการปั๊มพ่นสารชีวภัณฑ์แบบแปรผันตามโซน (Variable Rate Application: VRA)
\end{enumerate}
\end{tcolorbox}
